\documentclass[11pt]{article}
\usepackage{amsmath,amsthm,amssymb,mathtools,bbm}
\usepackage[margin=1in]{geometry}

\newtheorem{theorem}{Theorem}
\newtheorem{lemma}{Lemma}
\newtheorem{definition}{Definition}
\newtheorem{assumption}{Assumption}
\newtheorem{corollary}{Corollary}

\newcommand{\R}{\mathbb{R}}
\newcommand{\E}{\mathbb{E}}
\newcommand{\1}{\mathbbm{1}}
\newcommand{\ip}[2]{\left\langle #1,\, #2 \right\rangle}
\newcommand{\norm}[1]{\left\lVert #1 \right\rVert}
\newcommand{\grad}{\nabla}

\title{Importance-Weighted Randomized Coordinate Descent for Non-Convex, Coordinate-Wise Smooth Functions}
\author{}
\date{}

\begin{document}
\maketitle

\section*{Algorithm Name}
Importance-Weighted Randomized Coordinate Descent (IW-RCD).

\section*{Mathematical Formulation}
Let $f:\R^d\to\R$ be differentiable. Denote by $e_i\in\R^d$ the $i$-th standard basis vector and by $\nabla_i f(w)$ the $i$-th partial derivative at $w$. Let probabilities $p=(p_1,\dots,p_d)$ satisfy $p_i>0$ and $\sum_{i=1}^d p_i=1$.

At iteration $k\in\{0,1,2,\dots\}$:
- Sample a coordinate $i_k\in\{1,\dots,d\}$ with $\mathbb{P}(i_k=i)=p_i$.
- Update only coordinate $i_k$ according to
\[
w^{k+1} \;=\; w^k \;-\; \eta_k\,\frac{1}{p_{i_k}}\,\nabla_{i_k} f(w^k)\,e_{i_k},
\]
where $\eta_k>0$ is a stepsize.

This can be equivalently viewed as a stochastic-gradient step with the unbiased estimator
\[
g^k \;:=\; \frac{1}{p_{i_k}}\,\nabla_{i_k} f(w^k)\,e_{i_k},\qquad \E\!\left[g^k\mid w^k\right] \;=\; \nabla f(w^k).
\]

Stepsize constraint used in the analysis:
\[
0<\eta_k \;\le\; \min_{1\le i\le d}\frac{p_i}{L_i},
\]
where $L_i>0$ are coordinate-wise smoothness constants (defined below).

\section*{Pseudocode}
\noindent
Inputs: initial point $w^0\in\R^d$; probabilities $p\in\Delta^{d-1}$ with $p_i>0$; stepsizes $\{\eta_k\}_{k\ge 0}$; maximum iterations $T$ or tolerance $\varepsilon>0$.

\noindent
For $k=0,1,2,\dots$ until $k=T-1$ or $\norm{\nabla f(w^k)}\le \varepsilon$:
\begin{enumerate}
    \item Sample $i_k\sim p$.
    \item Compute the partial derivative $g_k:=\nabla_{i_k} f(w^k)$.
    \item Update: $w^{k+1} \leftarrow w^k - \eta_k \frac{1}{p_{i_k}} g_k\, e_{i_k}$.
\end{enumerate}
Return $w^T$ or the last iterate that satisfies the stopping criterion.

\section*{Dry Run Example}
Consider $d=1$ with $f(w)=(w-3)^2$, $w^0=0$, $\eta_k\equiv \eta=0.1$, and $p_1=1$.
- Gradient: $\nabla f(w)=2(w-3)$.
- Update: $w^{k+1}=w^k-\eta\cdot \nabla f(w^k)=w^k-0.2\,(w^k-3)=0.8\,w^k+0.6$.

Iterates and function values:
\begin{itemize}
\item $k=0$: $w^0=0$, $\nabla f(w^0)=2(0-3)=-6$, $w^1=0-0.1(-6)=0.6$, $f(w^0)=9$, $f(w^1)=(0.6-3)^2=5.76$.
\item $k=1$: $w^1=0.6$, $\nabla f(w^1)=2(0.6-3)=-4.8$, $w^2=0.6-0.1(-4.8)=1.08$, $f(w^2)=(1.08-3)^2=3.6864$.
\item $k=2$: $w^2=1.08$, $\nabla f(w^2)=2(1.08-3)=-3.84$, $w^3=1.08-0.1(-3.84)=1.464$, $f(w^3)=(1.464-3)^2=2.359296$.
\end{itemize}

\section*{Assumptions}
\begin{assumption}[Lower boundedness]
There exists $f_\star>-\infty$ such that $f(w)\ge f_\star$ for all $w\in\R^d$.
\end{assumption}

\begin{assumption}[Coordinate-wise smoothness]
For each coordinate $i\in\{1,\dots,d\}$, there exists $L_i>0$ such that for all $w\in\R^d$ and $h\in\R$,
\[
\left|\nabla_i f\!\left(w+h e_i\right)-\nabla_i f(w)\right| \le L_i |h|.
\]
Equivalently, for all $w\in\R^d$ and $h\in\R$,
\[
f(w+h e_i) \le f(w) + \nabla_i f(w)\,h + \frac{L_i}{2}h^2.
\]
\end{assumption}

\begin{assumption}[Sampling distribution]
Sampling probabilities $p=(p_1,\dots,p_d)$ satisfy $p_i>0$ and $\sum_{i=1}^d p_i=1$.
\end{assumption}

\begin{assumption}[Stepsizes]
The stepsizes $\{\eta_k\}_{k\ge 0}$ satisfy
\[
0<\eta_k \le \min_{1\le i\le d} \frac{p_i}{L_i}\qquad\text{for all }k\ge 0.
\]
\end{assumption}

\section*{Main Convergence Theorem}
\begin{theorem}[Expected descent and nonconvex stationarity rate]
\label{thm:main}
Under the assumptions above, the IW-RCD iterates $\{w^k\}_{k\ge 0}$ satisfy, for every $T\ge 1$,
\[
\sum_{k=0}^{T-1} \frac{\eta_k}{2}\, \E\big[\norm{\nabla f(w^k)}^2\big]
\;\le\; \E[f(w^0)] - f_\star \;=\; f(w^0)-f_\star.
\]
In particular, letting $\eta_{\min}:=\min_{0\le k\le T-1}\eta_k$ and choosing a random index $R$ uniformly in $\{0,\dots,T-1\}$ independent of the algorithm's history,
\[
\E\big[\norm{\nabla f(w^R)}^2\big] \;\le\; \frac{2\,(f(w^0)-f_\star)}{\eta_{\min}\,T}.
\]
If $\eta_k\equiv \eta \le \min_i p_i/L_i$, then $\E[\norm{\nabla f(w^R)}^2]=\mathcal{O}(1/T)$ with constant $2(f(w^0)-f_\star)/\eta$.
\end{theorem}

\section*{Theoretical Properties}
- Monotone expected descent: $\E[f(w^{k+1})]\le \E[f(w^k)]$ whenever $\eta_k \le 2\min_i p_i/L_i$. With the stricter choice $\eta_k \le \min_i p_i/L_i$, we obtain the uniform coefficient $1/2$ in Theorem~\ref{thm:main}.
- Importance sampling improves constants: Choosing $p_i \propto L_i$ maximizes $\min_i p_i/L_i$ subject to $\sum_i p_i=1$, yielding the largest allowable stepsize and best constants in the rate.
- Uniform sampling baseline: If $p_i=1/d$, then $\eta_k \le 1/(d\,L_{\max})$, where $L_{\max}:=\max_i L_i$, and the rate constant scales with $d\,L_{\max}$. With $p_i \propto L_i$, one can take $\eta_k \le 1/\sum_i L_i$, improving dependence from $d\,L_{\max}$ to $\sum_i L_i$.
- Hyperparameter sensitivity: The rate is inversely proportional to $\eta_{\min}$; picking $\eta_k$ close to $\min_i p_i/L_i$ improves the constant. Too large $\eta_k$ (exceeding $2\min_i p_i/L_i$) may destroy descent guarantees.

\section*{Discussion}
The algorithm behaves like stochastic gradient descent with an unbiased gradient estimator that updates a single coordinate per iteration. The importance factor $1/p_{i_k}$ is crucial to make the estimator unbiased and to achieve the descent bound that aggregates all coordinates symmetrically through $\norm{\nabla f(w^k)}^2=\sum_i \nabla_i f(w^k)^2$. Coordinate-wise smoothness suffices to control the one-dimensional local quadratic approximation along the sampled coordinate, yielding the standard nonconvex stationarity bound $\mathcal{O}(1/T)$ for the expected gradient norm squared.

\section*{Preliminaries (Definitions and Lemmas)}
We use the shorthand $\grad f(w)=(\nabla_1 f(w),\dots,\nabla_d f(w))^\top$ and $\norm{\cdot}$ for the Euclidean norm.

\begin{lemma}[Univariate descent along a coordinate]
\label{lem:coord_descent}
Under coordinate-wise smoothness, for any $w\in\R^d$, coordinate $i$, and scalar $d\in\R$,
\[
f(w+d\,e_i) \;\le\; f(w) + \nabla_i f(w)\,d + \frac{L_i}{2} d^2.
\]
Moreover, with $d=-\alpha\,\nabla_i f(w)$ for $\alpha\ge 0$,
\[
f(w-\alpha\,\nabla_i f(w)\,e_i) \;\le\; f(w) - \alpha\Bigl(1-\frac{L_i\alpha}{2}\Bigr)\,\nabla_i f(w)^2.
\]
\end{lemma}
\begin{proof}
Let $\phi(t):=f(w+t\,e_i)$ for $t\in\R$. Then $\phi'(t)=\nabla_i f(w+t\,e_i)$ and by coordinate-wise $L_i$-Lipschitz continuity of $\nabla_i f$ we have $|\phi'(t)-\phi'(0)|\le L_i |t|$. Integrating from $0$ to $d$,
\begin{align*}
\phi(d)-\phi(0)
&= \int_0^d \phi'(t)\,dt
= \int_0^d \bigl(\phi'(0) + (\phi'(t)-\phi'(0))\bigr)\,dt \\
&\le d\,\phi'(0) + \int_0^{|d|} L_i t \, dt
= d\,\nabla_i f(w) + \frac{L_i}{2} d^2,
\end{align*}
which yields the first inequality. Substituting $d=-\alpha\,\nabla_i f(w)$ gives
\[
f\bigl(w-\alpha\,\nabla_i f(w)\,e_i\bigr)
\le f(w) - \alpha\,\nabla_i f(w)^2 + \frac{L_i}{2}\alpha^2 \nabla_i f(w)^2
= f(w) - \alpha\Bigl(1-\frac{L_i\alpha}{2}\Bigr)\nabla_i f(w)^2.\qedhere
\]
\end{proof}

\begin{lemma}[One-step expected decrease]
\label{lem:one_step}
Let $w^{k+1}=w^k - \eta_k\,\frac{1}{p_{i_k}}\,\nabla_{i_k} f(w^k)\,e_{i_k}$ with $i_k\sim p$ independent of the past, and let $L_i$ be coordinate-wise smoothness constants. Then
\[
\E\!\left[f(w^{k+1}) \mid w^k\right]
\;\le\;
f(w^k) - \eta_k \sum_{i=1}^d \left(1 - \frac{L_i\eta_k}{2p_i}\right)\,\nabla_i f(w^k)^2.
\]
In particular, if $\eta_k \le \min_i p_i/L_i$, then
\[
\E\!\left[f(w^{k+1}) \mid w^k\right]
\;\le\;
f(w^k) - \frac{\eta_k}{2}\,\norm{\nabla f(w^k)}^2.
\]
\end{lemma}
\begin{proof}
Condition on $w^k$. For the sampled $i_k$, set $\alpha_{k,i_k}:=\eta_k/p_{i_k}\ge 0$. Apply Lemma~\ref{lem:coord_descent} with $i=i_k$ and $\alpha=\alpha_{k,i_k}$:
\begin{align*}
f(w^{k+1})
&= f\!\left(w^k - \alpha_{k,i_k}\,\nabla_{i_k} f(w^k)\,e_{i_k}\right)
\\
&\le f(w^k) - \alpha_{k,i_k}\Bigl(1-\tfrac{L_{i_k}\alpha_{k,i_k}}{2}\Bigr)\,\nabla_{i_k} f(w^k)^2
\qquad\textbf{[Step 1]: by Lemma \ref{lem:coord_descent} with } d=-\alpha_{k,i_k}\nabla_{i_k} f(w^k).
\end{align*}
Take conditional expectation over $i_k\sim p$:
\begin{align*}
\E\!\left[f(w^{k+1}) \mid w^k\right]
&\le f(w^k) - \sum_{i=1}^d p_i \left(\frac{\eta_k}{p_i}\right)\left(1-\frac{L_i}{2}\frac{\eta_k}{p_i}\right)\,\nabla_i f(w^k)^2
\qquad\textbf{[Step 2]: linearity of }\E\text{ and substitution } \alpha_{k,i}=\eta_k/p_i.
\\
&= f(w^k) - \eta_k \sum_{i=1}^d \left(1-\frac{L_i\eta_k}{2p_i}\right)\,\nabla_i f(w^k)^2
\qquad\textbf{[Step 3]: algebraic simplification}.
\end{align*}
If $\eta_k \le \min_i p_i/L_i$, then $0\le \frac{L_i\eta_k}{2p_i}\le \frac{1}{2}$ for all $i$, hence $1-\frac{L_i\eta_k}{2p_i}\ge \frac{1}{2}$. Therefore
\[
\E\!\left[f(w^{k+1}) \mid w^k\right]
\le f(w^k) - \frac{\eta_k}{2} \sum_{i=1}^d \nabla_i f(w^k)^2
= f(w^k) - \frac{\eta_k}{2}\,\norm{\nabla f(w^k)}^2
\qquad\textbf{[Step 4]: use }1-\tfrac{L_i\eta_k}{2p_i}\ge \tfrac{1}{2}\text{ and }\sum_i \nabla_i^2=\norm{\nabla}^2.
\]
\end{proof}

\section*{Main Proof (Step-by-Step Derivation)}
\begin{proof}[Proof of Theorem~\ref{thm:main}]
Start from Lemma~\ref{lem:one_step} and take total expectation using the tower property:
\begin{align*}
\E\big[f(w^{k+1})\big]
&= \E\big[\,\E[f(w^{k+1})\mid w^k]\,\big]
\qquad\textbf{[Step 5]: tower property } \E[X]=\E[\E[X\mid\mathcal{F}]].
\\
&\le \E\big[f(w^k) - \tfrac{\eta_k}{2}\,\norm{\nabla f(w^k)}^2\big]
\qquad\textbf{[Step 6]: apply Lemma \ref{lem:one_step} with }\eta_k \le \min_i p_i/L_i.
\\
&= \E[f(w^k)] - \frac{\eta_k}{2}\,\E\big[\norm{\nabla f(w^k)}^2\big]
\qquad\textbf{[Step 7]: linearity of expectation}.
\end{align*}
Rearrange:
\[
\frac{\eta_k}{2}\,\E\big[\norm{\nabla f(w^k)}^2\big]
\;\le\;
\E[f(w^k)] - \E[f(w^{k+1})]
\qquad\textbf{[Step 8]: move terms to the right-hand side}.
\]
Sum from $k=0$ to $T-1$:
\begin{align*}
\sum_{k=0}^{T-1} \frac{\eta_k}{2}\,\E\big[\norm{\nabla f(w^k)}^2\big]
&\le \sum_{k=0}^{T-1} \left(\E[f(w^k)] - \E[f(w^{k+1})]\right)
\qquad\textbf{[Step 9]: sum both sides over }k.
\\
&= \E[f(w^0)] - \E[f(w^T)]
\qquad\textbf{[Step 10]: telescoping sum } \sum_{k=0}^{T-1}(a_k-a_{k+1})=a_0-a_T.
\end{align*}
Use the lower bound $f(w^T)\ge f_\star$:
\[
\sum_{k=0}^{T-1} \frac{\eta_k}{2}\,\E\big[\norm{\nabla f(w^k)}^2\big]
\;\le\; f(w^0) - f_\star
\qquad\textbf{[Step 11]: apply Assumption 1}.
\]
This proves the first claim.

For the random index $R$ uniformly distributed in $\{0,\dots,T-1\}$ and independent of the history,
\begin{align*}
\E\big[\norm{\nabla f(w^R)}^2\big]
&= \frac{1}{T}\sum_{k=0}^{T-1} \E\big[\norm{\nabla f(w^k)}^2\big]
\qquad\textbf{[Step 12]: law of total expectation over discrete uniform }R.
\\
&\le \frac{2}{T}\sum_{k=0}^{T-1} \frac{1}{\eta_k}\,\left( \frac{\eta_k}{2}\E\big[\norm{\nabla f(w^k)}^2\big]\right)
\qquad\textbf{[Step 13]: multiply and divide by }\eta_k>0.
\\
&\le \frac{2}{T}\cdot \frac{1}{\eta_{\min}} \sum_{k=0}^{T-1} \left( \frac{\eta_k}{2}\E\big[\norm{\nabla f(w^k)}^2\big]\right)
\qquad\textbf{[Step 14]: since }\eta_k\ge \eta_{\min}.
\\
&\le \frac{2}{T}\cdot \frac{1}{\eta_{\min}} \,(f(w^0)-f_\star)
\qquad\textbf{[Step 15]: apply the first claim}.
\end{align*}
This concludes the proof.
\end{proof}

\section*{Rate Derivation (With Constants)}
- Uniform sampling $p_i=\tfrac{1}{d}$:
  - Stepsize constraint: $\eta \le \min_i \tfrac{1/d}{L_i} = \tfrac{1}{d\,L_{\max}}$ with $L_{\max}=\max_i L_i$.
  - Rate: for $\eta=\tfrac{1}{d\,L_{\max}}$, 
  \[
  \E\big[\norm{\nabla f(w^R)}^2\big] \le \frac{2d\,L_{\max}\,(f(w^0)-f_\star)}{T}.
  \]
- Importance sampling $p_i=\tfrac{L_i}{\sum_{j=1}^d L_j}$:
  - Stepsize constraint: $\eta \le \min_i \tfrac{L_i/\sum_j L_j}{L_i} = \tfrac{1}{\sum_j L_j}$.
  - Rate: for $\eta=\tfrac{1}{\sum_j L_j}$,
  \[
  \E\big[\norm{\nabla f(w^R)}^2\big] \le \frac{2\Bigl(\sum_{j=1}^d L_j\Bigr)\,(f(w^0)-f_\star)}{T}.
  \]
In general, the constant scales like $1/\eta_{\min}$, and larger $\min_i p_i/L_i$ allows larger $\eta$ and better constants.

\section*{Special Cases}
- One dimension ($d=1$): $p_1=1$, the update reduces to gradient descent with stepsize $\eta$, and Theorem~\ref{thm:main} recovers the standard nonconvex descent bound for $L_1$-smooth functions with $\eta\le 1/L_1$.
- Equal smoothness ($L_i\equiv L$): uniform sampling yields $\eta\le 1/(dL)$ and rate constant $2dL$; importance sampling coincides with uniform in this case.

\end{document}